\section{Bài tập chuyên đề}

\subsection{Củng cố}

\begin{tuluyen}
Hoàn thành các phương trình hóa học sau:
\begin{tasks}(2)
\task $\ce{Mg + O2 ->}$
\task $\ce{Al + O2 ->}$
\task $\ce{Fe + Cl2 ->}$
\task $\ce{Zn + HCl ->}$
\task $\ce{Fe + CuSO4 ->}$
\task $\ce{Cu + AgNO3 ->}$
\task $\ce{FeCl2 + NaOH ->}$
\task $\ce{FeCl3 + NaOH ->}$
\end{tasks}
\dapso{Viết và cân bằng các PTHH tương ứng.}
\end{tuluyen}

\begin{tuluyen}
Cho các cặp chất sau. Cặp nào có phản ứng trong dung dịch ở điều kiện thường?
\begin{tasks}(2)
\task Fe và \ce{CuSO4}.
\task Cu và \ce{HCl}.
\task Zn và \ce{AgNO3}.
\task Ag và \ce{ZnSO4}.
\end{tasks}
\dapso{Có phản ứng: Fe/\ce{CuSO4}, Zn/\ce{AgNO3}.}
\end{tuluyen}

\begin{tuluyen}
Cho \SI{5.60}{\gram} Fe tác dụng với \ce{HCl} dư. Tính thể tích \ce{H2} ở điều kiện có $V_m=\SI{24.8}{\liter\per\mole}$.
\dapso{\SI{2.48}{\liter}.}
\end{tuluyen}

\begin{tuluyen}
Cho \SI{13.0}{\gram} Zn vào \SI{200}{\milli\liter} dung dịch \ce{HCl} \SI{1.00}{\mole\per\liter}. Xác định chất hết, chất dư và thể tích \ce{H2} thu được.
\dapso{HCl hết; Zn dư \SI{0.100}{\mole}; $V_{\ce{H2}}=\SI{2.48}{\liter}$.}
\end{tuluyen}

\begin{tuluyen}
Nhúng Fe vào dung dịch \ce{CuSO4} dư. Khối lượng chất rắn tăng \SI{0.80}{\gram}. Tính khối lượng Cu bám trên thanh.
\dapso{\SI{6.40}{\gram}.}
\end{tuluyen}

\begin{tuluyen}
Hòa tan hoàn toàn \SI{9.00}{\gram} hỗn hợp Mg, Al trong \ce{HCl} dư, thu \SI{11.16}{\liter} \ce{H2}. Tính khối lượng mỗi kim loại.
\dapso{\SI{3.60}{\gram} Mg; \SI{5.40}{\gram} Al.}
\end{tuluyen}

\begin{tuluyen}
Một oxide sắt chứa $70.0\%$ Fe theo khối lượng. Xác định oxide trong ba chất \ce{FeO}, \ce{Fe2O3}, \ce{Fe3O4}.
\dapso{\ce{Fe2O3}.}
\end{tuluyen}

\begin{tuluyen}
Hoàn thành chuỗi phản ứng:
\[
\ce{Fe -> FeCl2 -> Fe(OH)2 -> Fe(OH)3 -> Fe2O3 -> Fe}.
\]
\dapso{Dùng \ce{HCl}, \ce{NaOH}, \ce{O2/H2O}, nhiệt, CO ở nhiệt độ cao.}
\end{tuluyen}

\begin{tuluyen}
Tính khối lượng Fe thu được khi khử hoàn toàn \SI{160}{\kilo\gram} \ce{Fe2O3} bằng CO.
\dapso{\SI{112}{\kilo\gram}.}
\end{tuluyen}

\begin{tuluyen}
Từ \SI{1.00}{\ton} quặng chứa $75\%$ \ce{Fe2O3}, hiệu suất thu hồi Fe là $90\%$. Tính khối lượng Fe thu được.
\dapso{\SI{0.4725}{\ton}.}
\end{tuluyen}

\begin{tuluyen}
Từ \SI{2.00}{\ton} quặng chứa $80\%$ \ce{ZnS}, phản ứng hoàn toàn theo sơ đồ \ce{ZnS -> ZnO -> Zn}. Tính khối lượng Zn.
\dapso{Khoảng \SI{1.072}{\ton}.}
\end{tuluyen}

\begin{tuluyen}
Nung \SI{2.00}{\ton} đá vôi chứa $85\%$ \ce{CaCO3}, hiệu suất phân hủy $80\%$. Tính khối lượng \ce{CaO} thu được.
\dapso{Khoảng \SI{0.762}{\ton}.}
\end{tuluyen}

\subsection{Bài tập tổng hợp}

\begin{tuluyen}
Trong công nghiệp, aluminium được sản xuất từ quặng bauxite bằng điện phân nóng chảy \ce{Al2O3}. Một mẫu quặng bauxite đã được làm giàu chứa $80\%$ \ce{Al2O3} về khối lượng, còn lại là tạp chất không chứa aluminium. Giả sử chỉ $85\%$ lượng \ce{Al2O3} trong quặng chuyển hóa thành kim loại. Để sản xuất \SI{500}{\kilo\metre} dây cáp nhôm, biết \SI{1}{\kilo\metre} dây cáp có khối lượng \SI{1080}{\kilo\gram}, tính khối lượng quặng bauxite cần dùng.
\source{Phỏng theo Câu 4d, đề tuyển sinh chuyên Hóa Lâm Đồng 2026.}
\dapso{\SI{1500}{\ton} quặng.}
\begin{loigiaiBT}
Khối lượng Al cần sản xuất:
\[
m_{\ce{Al}}=500\times1080=\SI{540000}{\kilo\gram}=\SI{540}{\ton}.
\]
\[
\ce{2Al2O3 -> 4Al + 3O2}\quad\text{(điện phân nóng chảy)}.
\]
Theo PTHH:
\[
m_{\ce{Al2O3},\text{lt}}=540\times\frac{102}{54}=\SI{1020}{\ton}.
\]
Do chỉ $85\%$ \ce{Al2O3} chuyển hóa:
\[
m_{\ce{Al2O3}}=\frac{1020}{0.85}=\SI{1200}{\ton}.
\]
Quặng chứa $80\%$ \ce{Al2O3}:
\[
m_{\text{quặng}}=\frac{1200}{0.80}=\SI{1500}{\ton}.
\]

\textbf{Vậy,} cần \SI{1500}{\ton} quặng bauxite.
\end{loigiaiBT}
\end{tuluyen}

\begin{tuluyen}
Tại một nhà máy luyện kim, quá trình tách Zn từ \SI{1.50}{\ton} quặng sphalerite chứa $80\%$ \ce{ZnS} được thực hiện qua hai giai đoạn:
\[
\ce{2ZnS + 3O2 ->[t^\circ] 2ZnO + 2SO2},
\]
\[
\ce{ZnO + C ->[t^\circ] Zn + CO}.
\]
Than cốc chứa $98\%$ C theo khối lượng. Các phản ứng tạo Zn xảy ra hoàn toàn.
\begin{enumerate}
\item Tính khối lượng Zn thu được và lượng than cốc cần dùng.
\item Toàn bộ \ce{SO2} được dùng sản xuất \ce{H2SO4} theo sơ đồ \ce{SO2 -> SO3 -> H2SO4}. Hiệu suất toàn quá trình là $80\%$. Tính thể tích dung dịch \ce{H2SO4} $98\%$ có khối lượng riêng \SI{1.84}{\gram\per\milli\liter} thu được.
\end{enumerate}
\source{Phỏng theo Câu 4.2, đề tuyển sinh chuyên Hóa Lâm Đồng 2026.}
\dapso{Zn $\approx\SI{0.804}{\ton}$; than cốc $\approx\SI{0.151}{\ton}$; dung dịch \ce{H2SO4} $\approx\SI{538}{\liter}$.}
\end{tuluyen}

\begin{tuluyen}
Hòa tan hỗn hợp Mg và Al trong acid dư. Khí \ce{H2} sinh ra được dẫn qua \ce{CuO} nung nóng; chỉ $80\%$ lượng \ce{H2} tham gia phản ứng. Sau phản ứng, khối lượng chất rắn giảm \SI{2.496}{\gram}. Tính số mol \ce{H2} đã sinh ra ban đầu.
\source{Phỏng theo kiểu dữ kiện gián tiếp của đề tuyển sinh chuyên Hóa Huế 2026.}
\dapso{$n_{\ce{H2}}=\SI{0.195}{\mole}$.}
\begin{loigiaiBT}
Khối lượng giảm chính là khối lượng oxygen bị lấy khỏi \ce{CuO}:
\[
n_{\ce{O}}=\frac{2.496}{16}=\SI{0.156}{\mole}.
\]
Theo \ce{CuO + H2 -> Cu + H2O}:
\[
n_{\ce{H2},\text{phản ứng}}=\SI{0.156}{\mole}.
\]
Do chỉ $80\%$ \ce{H2} tham gia phản ứng:
\[
n_{\ce{H2},\text{ban đầu}}=\frac{0.156}{0.80}=\SI{0.195}{\mole}.
\]

\textbf{Vậy,} ban đầu có \SI{0.195}{\mole} \ce{H2}.
\end{loigiaiBT}
\end{tuluyen}

\begin{tuluyen}
Hai loại quặng A, B chứa lần lượt $80\%$ và $60\%$ \ce{Fe2O3}. Người ta dùng tổng cộng \SI{500}{\ton} hai loại quặng để thu lượng \ce{Fe2O3} tương đương \SI{360}{\ton}. Tính khối lượng mỗi loại quặng.
\source{Phỏng theo kiểu phối trộn hai quặng trong đề tuyển sinh chuyên Hóa Thanh Hóa 2026.}
\dapso{A: \SI{300}{\ton}; B: \SI{200}{\ton}.}
\end{tuluyen}

\begin{tuluyen}
Trong lò cao, một phần oxide sắt được khử theo chuỗi \ce{Fe2O3 -> Fe3O4 -> FeO -> Fe}. Đá vôi được thêm vào để tạo \ce{CaO}, sau đó \ce{CaO} phản ứng với \ce{SiO2} tạo xỉ \ce{CaSiO3}.
\begin{enumerate}
\item Viết các phương trình hóa học của ba bước khử oxide sắt bằng CO.
\item Viết hai phương trình liên quan đến đá vôi và xỉ.
\item Giải thích vai trò của đá vôi trong quá trình luyện gang.
\end{enumerate}
\source{Phỏng theo bài luyện gang -- thép trong đề tuyển sinh chuyên Hóa Đắk Lắk 2026.}
\dapso{Xem phần lí thuyết về lò cao.}
\end{tuluyen}

\begin{tuluyen}
Một nhà máy nung \SI{100}{\ton} đá vôi chứa $92\%$ \ce{CaCO3}. Hiệu suất phân hủy đạt $90\%$. Khí \ce{CO2} sinh ra được thu hồi $75\%$. Tính khối lượng \ce{CaO} thu được và khối lượng \ce{CO2} được thu hồi.
\dapso{$m_{\ce{CaO}}\approx\SI{46.37}{\ton}$; $m_{\ce{CO2},\text{thu hồi}}\approx\SI{27.32}{\ton}$.}
\end{tuluyen}
